% !TeX encoding = UTF-8
% !TeX spellcheck = en_US
% !TeX TXS-program:compile = txs:///xelatex/[--shell-escape]

\documentclass[10pt]{article}

%%% FONT %%%
% \usepackage[utf8]{inputenc} % Remove inputenc when using XeLaTeX (fontspec handles UTF-8)
 \usepackage{fontspec}
 \setmainfont{Montserrat}[UprightFont = {Montserrat-Regular}, BoldFont = {Montserrat-Bold}]

\usepackage{graphicx} % gives \includegraphics
%\usepackage{ragged2e} % gives \justifying

\usepackage[usenames,dvipsnames]{color}

\usepackage{tikz} % this is to format the orange rectangles
\definecolor{mycolor}{HTML}{FA582D}
\definecolor{panorange}{HTML}{FF6A00}

%%% HEADER %%%
\usepackage{fancyhdr}
\pagestyle{fancy}
\fancyhf{}
\renewcommand{\headrulewidth}{0pt} % remove the underline from below the graphic
% \includegraphics[width=4cm]{../images/bio-header-palo.png}
\fancyhead[R]{\includegraphics[height=20pt]{../images/bio-header-palo.png}}

\setlength{\headheight}{25pt} % Adjust header height for graphics

\usepackage{setspace}
\onehalfspacing

\usepackage[hmargin=0.75in, vmargin=0.75in]{geometry}

\begin{document}

\vspace{0.3in}

\noindent\begin{tikzpicture}
	\fill[panorange] (0,0) rectangle (\linewidth,0.9cm);
	\node[text=white, font=\sffamily\bfseries] at ({0.5*\linewidth},0.45cm) {Franklin Diaz};
\end{tikzpicture}

\vspace{0.1in}

\noindent\fbox{\includegraphics[width=1.8in]{Franklin_24-09-19_3869.jpg}} \hspace{0.3in}
\begin{minipage}{\dimexpr\linewidth-1.8in-0.3in\relax}
\Large{\textbf{Principal Engineer, Security Researcher, and Educator}} \\[0.1in]
\small
Franklin Diaz is a Principal Engineer at Palo Alto Networks, where he specializes in public cloud infrastructure, automation, and DevSecOps as a member of the Professional Services team. In this role, he designs, develops, and deploys infrastructure as Code (IaC) for customers across public cloud environments (GCP, OCI, AWS, and Azure) based on best security practices and reference architectures, alongside extensive API automation work.

\vspace{0.15in}

A seasoned technology professional with over 20 years of experience, Diaz's background spans Telecommunications, Networks, Software, and Information Security. His career includes significant contributions such as his time at Salesforce as a Sr. Security Engineer, where his team designed, built, and maintained a large-scale security log collection solution processing over 1 billion security events per day. He also built and maintained data science pipelines, including one in the FedRamp High GovCloud environment in AWS, and developed a large-scale packet capture solution on internal company IT networks. Earlier in his career, as a Principal Research Engineer at Indiana University, he served as a Test Architect in the Software Defined Networking Lab, designing and directing test efforts on openflow networks.

\vspace{0.15in}

Diaz serves as an Adjunct Faculty member at the University of Colorado Denver Business School, teaching Information Systems courses and bridging the gap between industry practice and academic theory. His prior teaching experience includes five years at McHenry County College (2006--2011), where he taught Linux OS basics, Linux shell programming, Microsoft Office, and Microsoft Windows OS.

\vspace{0.15in}

Diaz holds an MS in Computer Information Systems from Northwestern University, an MS in Network Engineering/Network Security from DePaul University in Chicago, and a BS in Computer Science. He is an active contributor to the open-source community, notably as a contributor to Peirates, a Kubernetes penetration testing tool, and has run the BruCon 0x0F Capture the Flag contest.

\end{minipage}

\vspace{0.3in}

\noindent\begin{tikzpicture}
	\fill[panorange] (0,0) rectangle (\linewidth,0.75cm);
	\node[text=white, font=\sffamily\bfseries] at ({0.5*\linewidth},0.37cm) {Education};
\end{tikzpicture}

\vspace{0.05in}

\begin{itemize}[nosep]
	\item \textbf{MS, Computer Information Systems} --- Northwestern University
	\item \textbf{MS, Network Engineering/Network Security} --- DePaul University
	\item \textbf{BS, Computer Science} 
\end{itemize}

\vspace{0.15in}

\noindent\begin{tikzpicture}
	\fill[panorange] (0,0) rectangle (\linewidth,0.75cm);
	\node[text=white, font=\sffamily\bfseries] at ({0.5*\linewidth},0.37cm) {Teaching Experience};
\end{tikzpicture}

\vspace{0.05in}

\begin{itemize}[nosep]
	\item \textbf{University of Colorado Denver, Business School} --- Adjunct Faculty, Information Systems
	\item \textbf{McHenry County College} (2006--2011) --- Courses: Linux OS, Shell Programming, Microsoft Office, Windows OS
\end{itemize}

\vspace{0.3in}

\noindent\rule{\linewidth}{0.4pt}

\end{document}
